% GENERATED FILE. Edit canonical lessons, then run npm run generate:books.
% canonical-source-hash: fnv1a64:99e720e471e1059f
% canonical-lessons: UR-C523-istediyam, UR-C523-ampayar, UR-C523-iskor, UR-C523-nara, UR-C523-phandna

\chapter{At the Stadium}
\label{ch:ur-at-the-stadium}
\hlchaptermodality{523}

\begin{chapteropening}
\textbf{By the end of this chapter:} \emph{I can talk about the crowd and the score at a game.}

Complete the last lesson of chapter 523: I can talk about the crowd and the score at a game.
\end{chapteropening}

\section[\texorpdfstring{\ur{اسٹیڈیم} (isṭeḍiyam)}{اسٹیڈیم (isṭeḍiyam)}]{\ur{اسٹیڈیم} (isṭeḍiyam) — a stadium}

\label{lesson:UR-C523-istediyam}



\begin{quote}
Before the new one: say the Urdu for to make (someone) run, then the Urdu for a trophy.
\end{quote}

\subsection*{You'll want to know: \ur{اسٹیڈیم}}
\textbf{\ur{اسٹیڈیم}} — \emph{isṭeḍiyam} — \textquotedblleft{}a stadium\textquotedblright{}.

\begin{quote}
\ur{اسٹیڈیم} \ur{شائقین} \ur{سے} \ur{بھرا} \ur{تھا۔}
\end{quote}

\begin{quote}
\emph{isṭeḍiyam shāiqīn se bharā thā.}
\end{quote}

\begin{quote}
The stadium was full of fans.
\end{quote}

Masculine.

\begin{cousinweb}
\textbf{\ur{اسٹیڈیم}} (\emph{isṭeḍiyam}) is borrowed from English \emph{stadium}.
\end{cousinweb}

\begin{grammarlens}[title={What you've built}]
One more word for describing what you lived and felt.
\end{grammarlens}

\subsection*{Your turn}
\begin{itemize}
\raggedright
  \item \emph{Say it:} \emph{isṭeḍiyam}
  \item \emph{Say it:} \emph{isṭeḍiyam}, once more
  \item \emph{Say it:} \emph{isṭeḍiyam shāiqīn se bharā thā.} — The stadium was full of fans.
  \item \emph{Recall:} say the Urdu for to make (someone) run, then the Urdu for a trophy, then say \emph{isṭeḍiyam} again
\end{itemize}

\begin{tcolorbox}[breakable,colback=teal!4,colframe=teal!35!black,title={Before you move on}]
What does \ur{اسٹیڈیم} mean? (\textquotedblleft{}A stadium\textquotedblright{}.) Say it once more.
\end{tcolorbox}
\section[\texorpdfstring{\ur{امپائر} (ampāyar)}{امپائر (ampāyar)}]{\ur{امپائر} (ampāyar) — an umpire}

\label{lesson:UR-C523-ampayar}



\begin{quote}
Before the new one: say the Urdu for a trophy, then the Urdu for a stadium.
\end{quote}

\subsection*{You'll want to know: \ur{امپائر}}
\textbf{\ur{امپائر}} — \emph{ampāyar} — \textquotedblleft{}an umpire\textquotedblright{}.

\begin{quote}
\ur{امپائر} \ur{نے} \ur{کھیل} \ur{روکا۔}
\end{quote}

\begin{quote}
\emph{ampāyar ne khel rokā.}
\end{quote}

\begin{quote}
The umpire stopped the game.
\end{quote}

Masculine or feminine, by the person meant.

\begin{cousinweb}
\textbf{\ur{امپائر}} (\emph{ampāyar}) is borrowed from English \emph{umpire}.
\end{cousinweb}

\begin{grammarlens}[title={What you've built}]
One more word for describing what you lived and felt.
\end{grammarlens}

\subsection*{Your turn}
\begin{itemize}
\raggedright
  \item \emph{Say it:} \emph{ampāyar}
  \item \emph{Say it:} \emph{ampāyar}, once more
  \item \emph{Say it:} \emph{ampāyar ne khel rokā.} — The umpire stopped the game.
  \item \emph{Recall:} say the Urdu for a trophy, then the Urdu for a stadium, then say \emph{ampāyar} again
\end{itemize}

\begin{tcolorbox}[breakable,colback=teal!4,colframe=teal!35!black,title={Before you move on}]
What does \ur{امپائر} mean? (\textquotedblleft{}An umpire\textquotedblright{}.) Say it once more.
\end{tcolorbox}
\section[\texorpdfstring{\ur{اسکور} (iskor)}{اسکور (iskor)}]{\ur{اسکور} (iskor) — a score (in a game)}

\label{lesson:UR-C523-iskor}



\begin{quote}
Before the new one: say the Urdu for a stadium, then the Urdu for an umpire.
\end{quote}

\subsection*{You'll want to know: \ur{اسکور}}
\textbf{\ur{اسکور}} — \emph{iskor} — \textquotedblleft{}a score (in a game)\textquotedblright{}.

\begin{quote}
\ur{میچ} \ur{کا} \ur{اسکور} \ur{کیا} \ur{ہے؟}
\end{quote}

\begin{quote}
\emph{maich kā iskor kyā hai?}
\end{quote}

\begin{quote}
What's the score in the match?
\end{quote}

Masculine.

\begin{cousinweb}
\textbf{\ur{اسکور}} (\emph{iskor}) is borrowed from English \emph{score}.
\end{cousinweb}

\begin{grammarlens}[title={What you've built}]
One more word for describing what you lived and felt.
\end{grammarlens}

\subsection*{Your turn}
\begin{itemize}
\raggedright
  \item \emph{Say it:} \emph{iskor}
  \item \emph{Say it:} \emph{iskor}, once more
  \item \emph{Say it:} \emph{maich kā iskor kyā hai?} — What's the score in the match?
  \item \emph{Recall:} say the Urdu for a stadium, then the Urdu for an umpire, then say \emph{iskor} again
\end{itemize}

\begin{tcolorbox}[breakable,colback=teal!4,colframe=teal!35!black,title={Before you move on}]
What does \ur{اسکور} mean? (\textquotedblleft{}A score (in a game)\textquotedblright{}.) Say it once more.
\end{tcolorbox}
\section[\texorpdfstring{\ur{نعرہ} (nārā)}{نعرہ (nārā)}]{\ur{نعرہ} (nārā) — a cheer; a slogan}

\label{lesson:UR-C523-nara}



\begin{quote}
Before the new one: say the Urdu for an umpire, then the Urdu for a score.
\end{quote}

\subsection*{You'll want to know: \ur{نعرہ}}
\textbf{\ur{نعرہ}} — \emph{nārā} — \textquotedblleft{}a cheer; a slogan\textquotedblright{}.

\begin{quote}
\ur{اسٹیڈیم} \ur{میں} \ur{ہر} \ur{طرف} \ur{نعروں} \ur{کی} \ur{آواز} \ur{تھی۔}
\end{quote}

\begin{quote}
\emph{isṭeḍiyam meṅ har taraf nāroṅ kī āvāz thī.}
\end{quote}

\begin{quote}
The sound of cheering filled the stadium.
\end{quote}

Masculine.

\begin{cousinweb}
\textbf{\ur{نعرہ}} (\emph{nārā}) is usually traced, through Persian, to Arabic \emph{naʿra}, \textquotedblleft{}a shout, a cry\textquotedblright{}.
\end{cousinweb}

\begin{grammarlens}[title={What you've built}]
One more word for describing what you lived and felt.
\end{grammarlens}

\subsection*{Your turn}
\begin{itemize}
\raggedright
  \item \emph{Say it:} \emph{nārā}
  \item \emph{Say it:} \emph{nārā}, once more
  \item \emph{Say it:} \emph{isṭeḍiyam meṅ har taraf nāroṅ kī āvāz thī.} — The sound of cheering filled the stadium.
  \item \emph{Recall:} say the Urdu for an umpire, then the Urdu for a score, then say \emph{nārā} again
\end{itemize}

\begin{tcolorbox}[breakable,colback=teal!4,colframe=teal!35!black,title={Before you move on}]
What does \ur{نعرہ} mean? (\textquotedblleft{}A cheer; a slogan\textquotedblright{}.) Say it once more.
\end{tcolorbox}
\section[\texorpdfstring{\ur{پھاندنا} (phāṅdnā)}{پھاندنا (phāṅdnā)}]{\ur{پھاندنا} (phāṅdnā) — to leap over}

\label{lesson:UR-C523-phandna}



\begin{quote}
Before the new one: say the Urdu for a score, then the Urdu for a cheer.
\end{quote}

\subsection*{You'll want to know: \ur{پھاندنا}}
\textbf{\ur{پھاندنا}} — \emph{phāṅdnā} — \textquotedblleft{}to leap over\textquotedblright{}.

\begin{quote}
\ur{گھوڑے} \ur{نے} \ur{دیوار} \ur{پھاندی۔}
\end{quote}

\begin{quote}
\emph{ghoṛe ne dīvār phāṅdī.}
\end{quote}

\begin{quote}
The horse leapt over the wall.
\end{quote}

Past, with \ur{نے} (he or she …): \ur{اس} \ur{نے} \ur{دیوار} \ur{پھاندی}, \emph{us ne dīvār phāṅdī}.

\begin{cousinweb}
\textbf{\ur{پھاندنا}} (\emph{phāṅdnā}) is inherited: CDIAL lists the Hindi cousin \emph{phā̃dnā}, \textquotedblleft{}to leap\textquotedblright{}, under Sanskrit \emph{spandatē}, \textquotedblleft{}trembles, kicks\textquotedblright{} (CDIAL 13806).
\end{cousinweb}

\begin{grammarlens}[title={What you've built}]
One more word for describing what you lived and felt.
\end{grammarlens}

\subsection*{Your turn}
\begin{itemize}
\raggedright
  \item \emph{Say it:} \emph{phāṅdnā}
  \item \emph{Say it:} \emph{phāṅdnā}, once more
  \item \emph{Say it:} \emph{ghoṛe ne dīvār phāṅdī.} — The horse leapt over the wall.
  \item \emph{Recall:} say the Urdu for a score, then the Urdu for a cheer, then say \emph{phāṅdnā} again
\end{itemize}

\begin{tcolorbox}[breakable,colback=teal!4,colframe=teal!35!black,title={Before you move on}]
What does \ur{پھاندنا} mean? (\textquotedblleft{}To leap over\textquotedblright{}.) Say it once more.
\end{tcolorbox}
